%!TEX root =  main.tex


\para{Irrecoverable from biased batch workload.}
The workload balance between batches can become irrecoverably disrupted under certain conditions. When a request is terminated due to user abort or preemption, \alg{} may fail to recover from the resulting imbalance because new replacement running requests are consistently assigned to the draft batch. While this policy ensures that draft tokens are available for verification, it prevents the load-balancing mechanism (the \texttt{assign} operation) from correcting the induced batch imbalance.
Furthermore, when chunked prefill is enabled, the Scheduler may initialize fewer than $2m$ running requests in the first SD step. In this case, \alg{} must continue to apply the load-balancing \texttt{assign} operation for all subsequent batch ID assignments, rather than switching to the simplified post–first-step policy. Otherwise, an irrecoverable workload imbalance between the two batches can arise.
However, continuously applying the \texttt{assign} operation introduces a minor overhead. If a new running request is assigned to the target batch, the Verifier can sample only a single token because there are no draft tokens, resulting in one less efficient “idle” SD step for that request. Since this cost is incurred at most once per request, it is deemed an acceptable trade-off for the long-term performance benefits of maintaining batch balance. Accordingly, we plan to adopt this revised batch ID assignment strategy in a future version of \alg{} that supports the vLLM v1 engine and chunked prefill.


\para{Batch exhaustion and tail effect.}
The Batch Parallelism scheme, as specified in \Cref{sec:minedraft}, includes a fallback mechanism: when one batch becomes empty and the other batch contains no verifiable drafts, the system reverts to standard SD, which inevitably curtails the duration of parallel execution. Although Batch Parallelism successfully parallelizes the substantial majority of SD steps and yields significant performance gains per our empirical results, the tail effect circumscribes its maximum potential improvement.
A possible approach to offset this limitation involves adopting a modified PEARL parallelism. When one batch becomes empty, the system instructs the Drafter to generate draft tokens for the remaining batch, including those undergoing verification, without confirming their success. In the next step, instead of verifying only previously successful requests as in standard PEARL, the system re-drafts requests that failed verification before submitting the batch to the Verifier. This re-drafting mechanism for failed verification attempts is intended for partially overlapping drafting time, optimizing performance to the greatest extent possible.
